% titel: Blatt 6, A6: Polynome 3-ten Grades (Cardano)
% skript: Kap. 1 (Wurzeln), Kap. 3 (Nullstellen)

% Teil I: Werkzeuge, kein Lösungsweg.

Mit der Cardano-Formel eine reelle Nullstelle kubischer Polynome der Form
$z^3+3pz+2q$ konstruieren und alle Nullstellen eines konkreten Polynoms
bestimmen.

\begin{werkzeugkasten}[Die Grundidee hinter Cardano -- ein Trick mit dem Binom]
Für eine quadratische Gleichung gab es eine fertige Lösungsformel. Für die
kubische Gleichung $z^3+3pz+2q=0$ (schon von der $z^2$-Form befreit, was
so eine Verschiebung der Unbestimmten bewirkt) rät man eine Struktur: Eine
Summe $x=\sqrt[3]{u}+\sqrt[3]{v}$ mit \emph{passend} gewählten $u,v$.
Warum gerade diese? Weil das Binom $(a+b)^3=a^3+b^3+3ab(a+b)$ die Summe
$(\sqrt[3]{u}+\sqrt[3]{v})^3$ so zurückführt, dass genau der lineare Term
$3ab\,x$ entsteht -- und der Term $3px$ in der Gleichung genau dazu passt.

Konkret: setzt man $a=\sqrt[3]{u},\ b=\sqrt[3]{v}$, dann ist
\begin{equation*}
  x^3=u+v+3\underbrace{\sqrt[3]{uv}}_{\mathstrut -p}\,x .
\end{equation*}
Ordnet man $u+v=-2q$ und $\sqrt[3]{uv}=-p$, so ist $x$ automatisch eine
Lösung. Die Cardano-Wahlen
\begin{equation*}
  u=-q+\sqrt{D},\qquad v=-q-\sqrt{D},\qquad D:=q^2+p^3
\end{equation*}
machen genau das: $u+v=-2q$, und $uv=q^2-D=-p^3$, also
$\sqrt[3]{uv}=-p$. Der wichtige Einzelschritt ist, zuerst
$\sqrt[3]{uv}=-p$ zu berechnen -- er \emph{entschärft} den $3px$-Term.

\textit{Skript:} Kap. 1 (komplexe Wurzeln), Kap. 3 (Polynome).
\end{werkzeugkasten}

\begin{werkzeugkasten}[Die reelle dritte Wurzel]
\emph{$\sqrt[3]{\cdot}$} meint hier die \emph{reelle} Kubikwurzel: die
Umkehrung von $t\mapsto t^3$. Sie existiert immer, auch für negative
Argumente ($\sqrt[3]{-9}=-\sqrt[3]{9}$), denn jede reelle Zahl hat genau
eine reelle dritte Potenz/Wurzel. Für die Cardano-Voraussetzung
$D=q^2+p^3>0$ bleibt man im Reellen; der Fall $D<0$ (drei reelle
Nullstellen, \emph{casus irreducibilis}) führt über die komplexen Zahlen.

\textit{Skript:} Kap. 1 (Wurzeln).
\end{werkzeugkasten}

\begin{werkzeugkasten}[Alle Nullstellen durch Faktorisieren]
Hat man eine reelle Nullstelle $x$, faktorisiert man das kubische Polynom:
\begin{equation*}
  z^3+3pz+2q=(z-x)\bigl(z^2+xz+(x^2+3p)\bigr).
\end{equation*}
Die Verifikation durch Ausmultiplizieren: $(z-x)(z^2+xz+(x^2+3p))$ liefert
$z^3+3pz-x^3-3px$; wegen $x^3+3px=-2q$ ist das $z^3+3pz+2q$. Die
quadratische Gleichung $z^2+xz+(x^2+3p)=0$ liefert die restlichen zwei
(ggf. konjugiert komplexen) Nullstellen.

\textit{Skript:} Kap. 3 (Nullstellen von Polynomen).
\end{werkzeugkasten}

\begin{achtung}[Typische Fallen]
\begin{itemize}
  \item \textbf{Reelle Kubikwurzel nicht verwechseln} mit der (mehrdeutigen)
        komplexen: $\sqrt[3]{-9}=-(\sqrt[3]{9})$, nicht $2{,}080\,\mathrm e^{i\pi/3}$.
  \item \textbf{Zuerst $\sqrt[3]{uv}$ berechnen:} Der schnelle Weg nutzt
        $uv=-p^3$, nicht ein Ausmultiplizieren von $(a+b)^3$.
  \item \textbf{Beim konkreten Polynom $p$ und $q$ erkennen:}
        $3z\sqrt[3]{9}$ heisst $3p=3\sqrt[3]{9}$ und $2q=8$; nicht $p$ mit
        $3p$ verwechseln.
  \item \textbf{Alle drei Nullstellen angeben}, nicht nur die reelle.
\end{itemize}
\end{achtung}
